\section{Classical Field Theory as a Bridge Between Classical Mechanics and Quantum Field Theory}

Classical field theory is one of the central languages of modern theoretical
physics. It begins from ideas that are already familiar in classical mechanics:
configuration, motion, action, symmetry, and conservation. It then replaces a
finite list of generalized coordinates by fields, whose degrees of freedom are
spread continuously over space and time.

In mechanics, a system may be described by coordinates
\[
    q^1(t),\ldots,q^n(t).
\]
In field theory, the dynamical variables are functions such as
\[
    \phi(t,\vect{x}), \qquad A_\mu(x), \qquad \psi_\alpha(x).
\]
The difference is not merely notational. A field has independent values at
different points, and its equations usually involve partial derivatives with
respect to both time and position. This is why field theory is the natural
language of waves, electromagnetism, continua, gauge interactions, and the
classical limit of quantum field theory.

These notes are written as a bridge. On one side stands classical mechanics,
especially the Lagrangian and Hamiltonian viewpoints. On the other side stand
quantum field theory, gauge theory, and relativistic physics. The aim is not to
quantize fields in this course, but to make the classical structures precise
enough that their later quantization is not mysterious.

\begin{remarkbox}
Classical field theory is not a historical leftover. It is the structural
language behind many modern theories. Even when a theory is ultimately quantum,
its action, symmetries, currents, constraints, and field equations are usually
first understood classically.
\end{remarkbox}
